\label{chp:sr12}

In this section, we will examine the electromagnetic form factors of $\nicefrac{3}{2}\rightarrow\nicefrac{1}{2}$ transitions between the double heavy baryons with quark content $QQ'q$.This transition in the presence of the electromagnetic field is described by
\begin{align}
    \Pi^{\mu\nu}(p,q) = i\int d^4x \int d^4y e^{ipx+iqy}\mel{0}{\mathcal{T}\left\{\eta_{\nicefrac{1}{2}}(x)j_{el}^\nu(y)\overline{\eta}^\mu_{\nicefrac{3}{2}}(0)\right\}}{0}
    \label{initial_corr}
\end{align}
%
where $j_{el}^\nu$ is the electromagnetic current given by
%
\begin{equation}
    j_{el}^\nu = \sum_{q'} e_{q'} \bar q' \gamma^\nu q'
\end{equation}
where $e_{q'}$ is the charge of the quark $q'$ and the sum over is over all quark flavors. Note that although the electromagnetic current contains a sum over all quarks, only the terms proportional to the charges of the quarks appearing in the double heavy baryon contribute to the transition at the order studied in this work. Here $\eta_{\nicefrac{1}{2}}$ and $\eta^\mu_{\nicefrac{3}{2}}$ are the interpolating currents of spin-$\nicefrac{1}{2}$ and spin-$\nicefrac{3}{2}$ baryons respectively.

The $J^P$ quantum numbers of doubly heavy baryons are given in Table-\ref{baryontable}. Note that $\Xi^\prime_{QQ^\prime q}$ is antisymmetric under the exchange of heavy quarks whereas the remaining states are symmetric in the flavor space.

\begin{table}[h!]
    \centering
    \begin{tabular}{|c|c|}
        \hline
        State & J$^\text{P}$  \\
        \hline
        $\Xi_{QQq}$&$\nicefrac{1}{2}^+$\\
        \hline
        $\Xi_{QQq}^\ast$&$\nicefrac{3}{2}^+$\\
        \hline
        $\Xi_{QQ^\prime q}$&$\nicefrac{1}{2}^+$\\
        \hline
        $\Xi_{QQ^\prime q}^\ast$&$\nicefrac{3}{2}^+$\\
        \hline
        $\Xi^\prime_{QQ^\prime q}$&$\nicefrac{1}{2}^+$\\
        \hline
    \end{tabular}
    \caption{$J^P$ quantum numbers of doubly heavy $\Xi$ baryons}
    \label{baryontable}
\end{table}

\begin{equation}
    \begin{aligned}
        \eta^{(S)} &= \dfrac{N}{2}\epsilon^{abc}\sum_{k=1}^2\left[\left(Q^{aT}CA_1^kq^b\right)A_2^kQ^{\prime c}+\left(Q^{\prime aT}CA_1^kq^b\right)A_2^kQ^c\right] \\
        \eta^{(A)} &= \dfrac{1}{\sqrt{6}}\epsilon^{abc}\sum_{k=1}^2\left[2\left(Q^{aT}CA_1^kQ^{\prime b}\right)A_2^kq^c + \left(Q^{aT}CA_1^kq^b\right)A_2^kQ^{\prime c} - \left(Q^{\prime aT}CA_1^kq^b\right)A_2^kQ^{c}\right]
    \end{aligned}
    \label{spin12currents}
\end{equation}
%
where $a,b,c$ are color indices, $N$ is the normalization constant which is equal to $\sqrt{2}$ when $Q\neq Q^\prime$, $1$ otherwise, $C$ is the charge conjugation matrix, $T$ is the transposition, $A_1^1 = 1$, $A_1^2 = \beta\gamma_5$, $A_2^1 = \gamma_5$ and $A_2^2 = \beta$, and $\beta$ is an arbitrary parameter. Superscripts $(S)$ and $(A)$ mean symmetric and anti-symmetric interpolating currents with respect to the exchange of heavy quarks.

The interpolating current of doubly heavy baryon with spin-$\nicefrac{3}{2}$ can be written as

\begin{equation}
    \eta^{\mu} = \dfrac{N}{\sqrt{3}}\varepsilon^{abc}\left[\left(q^{aT}C\gamma^{\mu}Q^b\right)Q^{\prime c} + \left(q^{aT}C\gamma^{\mu}Q^{\prime b}\right)Q^c + \left(Q^{aT}C\gamma^\mu Q^{\prime b}\right)q^c\right]
    \label{spin32current}
\end{equation}
where again $a,b,c$ are color indices, $N$ is the normalization constant which is equal to $\sqrt{2}$ when $Q\neq Q^\prime$, $1$ otherwise, $C$ is the charge conjugation operator. Note that this interpolating current is symmetric under the exchange of any two quark fields.

By introducing the plane wave electromagnetic background field

\begin{equation}
    F_{\mu\nu} = i\left(\varepsilon_\nu^{(\lambda)}q_\mu-\varepsilon_\mu^{(\lambda)}q_\nu\right)e^{iqx}
    \label{backgroundfield}
\end{equation}
and multiplying the correlation function (\ref{initial_corr}) by $\varepsilon_\nu^{(\lambda)}$, it can be rewritten in the following form
\begin{equation}
    \Pi^{\mu\nu}(p,s) \varepsilon_\nu^{(\lambda)} = i\int d^4x e^{ipx} \mel{0}{\mathcal{T}\left\{\eta_{\nicefrac{1}{2}}(x)\overline{\eta}^\mu_{\nicefrac{3}{2}}(0)\right\}}{0}_F
    \label{spin12bfcorr}
\end{equation}

In Eq-(\ref{spin12bfcorr}) the subscript $F$ indicates that all vacuum expectation values are evaluated in the presence of the background field. 
It should be noted that the correlation function given in Eq-(\ref{initial_corr}) can be obtained from Eq-(\ref{spin12bfcorr}) by expanding latter in powers of the background field and taking into account only terms linear in $F_{\mu\nu}$, which corresponds to the single photon emission. (More about the background field method can be found in \cite{PhotonDA} and \cite{ExternalField}.)

Now let us derive the sum rules for the $\nicefrac{3}{2}\longrightarrow\nicefrac{1}{2}$ transition form factors.  
According to the sum rules philosophy, the correlation function Eq-(\ref{initial_corr}) has to be calculated in two different kinematic domains. 
The hadronic representation of the correlation function can be calculated by inserting a full set of hadrons carrying the same quantum numbers as the interpolating currents $\eta_{\nicefrac{1}{2}}$ and $\eta^\mu_{\nicefrac{3}{2}}$ and 
then separating the contributions of ground states. On the other hand, the correlation function can be calculated in deep Euclidean region, where $p^2\ll0$ and $(p+q)^2\ll0$ with the help of operator product expansion (OPE). 
The OPE is performed in terms of photon distribution amplitudes(DAs) of increasing twist. 

At the hadronic level, the correlation function Eq-(\ref{spin12bfcorr}) can be written as
\begin{equation}
    \Pi^{\mu\nu}\epsilon_\nu^{(\lambda)} = \dfrac{\epsilon_\nu^{(\lambda)}\mel{0}{\eta_{\nicefrac{1}{2}}}{B_2(p,s)}\mel{B_2(p,s)}{j_{el}^\nu}{B_1(p+q,s^\prime}\mel{B_1(p+q,s^\prime)}{\overline{\eta}^\mu_{\nicefrac{3}{2}}}{0}}{\left(p^2-m_2^2\right)\left((p+q)^2-m_1^2\right)}
    \label{corr_with_inserted_states}
\end{equation}
where a sum over hadrons and the continuum is not shown explicitly.
The matrix elements in eq-(\ref{corr_with_inserted_states}) are defined as
\begin{equation}
    \begin{aligned}
        \mel{0}{\eta_{\nicefrac{1}{2}}}{B_{\nicefrac{1}{2}}(p,s)} &= \lambda_2 u(p,s)\\
        \mel{B_{\nicefrac{3}{2}}(p+q,s)}{\overline{\eta}^\mu_{\nicefrac{3}{2}}}{0} &= \lambda_1 \overline{u}^\mu(p+q,s)
    \end{aligned}
    \label{mels}
\end{equation}
where $u$ is a Dirac spinor, $u^\mu$ is a Rarita-Schwinger spinor and $\lambda_{1\left(2\right)}$ is the residue of the spin-$\nicefrac{3}{2}\left(\nicefrac{1}{2}\right)$ doubly heavy baryon.

The matrix element of the  electromagnetic current sandwiched between spin-$\nicefrac{1}{2}$ and spin-$\nicefrac{3}{2}$ states is parametrized in terms of three form factors in the following way \cite{JONESmatrixel,KORNERmatrixel}:
\begin{equation}
    \begin{aligned}
        \mel{B_{\nicefrac{1}{2}}(p,s)}{j_{el}^\nu}{B_{\nicefrac{3}{2}}(p+q,s)} &=\overline{u}(p,s)\Bigg[G_1\left(-g^{\alpha\nu}\slashed{q}+q^{\alpha}\gamma^{\nu}\right)\\
        &+G_2\left(-g^{\alpha\nu}\left(p+\dfrac{q}{2}\right)q + q^\alpha\left(p+\dfrac{q}{2}\right)^\nu\right)\\
        &+G_3\left(q^\alpha q^\nu - q^2g^{\alpha\nu}\right)\Bigg]\gf u_\alpha(p+q,s)
    \end{aligned}
    \label{el_mel}
\end{equation}
The last term which is proportional to $ G_3$ does not contribute to the emission of a real photon. In the problem studied in this work, the emitted photon is a real photon. For this reason, only the values of $G_1$ and $G_2$ at $q^2=0$ are needed.

Magnetic dipole ($G_E$) and electric quadrupole ($G_E$) moments, which are more directly accessible in experiments, are defined through the formfactors $G_1$ and $G_2$ as
\cite{JONESmatrixel,KORNERmatrixel}:
%
\begin{equation}
    \begin{aligned}
        G_M &= \Bigg[\dfrac{3m_1+m_2}{m_1}G_1(0) + \left(m_1-m_2\right)G_2(0)\Bigg]\dfrac{m_2}{3}\\
        G_E &= \left(m_1-m_2\right)\Bigg[\dfrac{G_1(0)}{m_1}+G_2(0)\Bigg]\dfrac{m_2}{3}
    \end{aligned}
    \label{GMandGE}
\end{equation}
%
where $m_1$ and $m_2$ are masses of spin-$\nicefrac{3}{2}$ and spin-$\nicefrac{1}{2}$ doubly heavy baryons, respectively.
%
Performing summation over spins for the Dirac and Rarita-Schwinger spinors via
\begin{equation}
    \begin{aligned}
        \sum u(p,s)\overline{u}(p,s) &= \left(\slashed{p}+m\right)\\
        \sum u_\alpha(p,s)\overline{u}_\beta(p,s) &= \left(\slashed{p}+m\right)\left[-g_{\alpha\beta}+\dfrac{1}{3}\gamma_\alpha\gamma_\beta-\dfrac{2p_\alpha p_\beta}{3m^2}-\dfrac{p_\alpha\gamma_\beta-p_\beta\gamma_\alpha}{3m}\right]
    \end{aligned}
    \label{spin_summations}
\end{equation}
the phenomenological representation of the correlation function can be obtained. In order to obtain the sum rules for transition form factors, two issues need to be addressed:
\begin{enumerate}
    \item Besides the spin-$\nicefrac{3}{2}$ baryon, the current $\eta_\mu$ can also create a negative parity baryon with spin-$\nicefrac{1}{2}$ from the vacuum. This negative parity baryon also contributed to the correlation function.
    \item Not all Lorentz structures are independent of each other.
\end{enumerate}
The first problem can be solved in the following way. The matrix element of the current $\eta_\mu$ between vacuum and a spin-$\nicefrac{1}{2}$ baryon can be written as
%
\begin{equation}
    \mel{B_{\nicefrac{1}{2}^-}(p+q,s)}{\eta_\mu}{0} = A \bar u(p+q) \left[(p+q)_\mu-\dfrac{m}{4}\gamma_\mu\right]
\end{equation}
where the fact that $\eta_\mu \gamma^\mu=0$ is used. Hence, it is seen that the contributions of the spin-$\nicefrac{1}{2}$ baryons are either $\propto (p+q)_\mu$ or have $\gamma_\mu$ at the right. Other structures do not receive any contribution from these spin-$\nicefrac{1}{2}$ baryons, and only receive contribution from spin-$\nicefrac{3}{2}$ baryons.
%
In order to solve the second issue and obtain independent Lorentz structures (following \cite{AlievRadDec}) the Dirac matrices are ordered as $\slashed{\varepsilon}\slashed{q}\slashed{p}\gamma_\mu$. With this ordering, the hadronic representation of the correlation function becomes:
%
    \begin{align}
        \nonumber\Pi_{\mu\nu}(p,q)\varepsilon^\nu =& e \lambda_2 \lambda_1 \frac{1}{p^{2}-m_2^{2}} \frac{1}{(p+q)^{2}-m_1^{2}}\Bigg\{\\
        \nonumber&{\left[\varepsilon_{\mu}(p q)-(\varepsilon p) q_{\mu}\right]\left\{-2 G_{1} m_1-G_{2} m_1 m_2+G_{2}(p+q)^{2}\right.}\\
        \nonumber+&\left.\left[2 G_{1}+G_{2}\left(m_2-m_1\right)\right] \slashed{p}+m_2 G_{2} \slashed{q}-G_{2} \slashed{q} \slashed{p}\right\} \gamma_{5}\\
        +&\left[q_{\mu} \slashed{\varepsilon}-\varepsilon_{\mu} \slashed{q}\right]\left\{G_{1}\left(p^{2}+m_1 m_2\right)-G_{1}\left(m_1+m_2\right) \slashed{p}\right\} \gamma_{5}\\
        \nonumber+&2 G_{1}[\slashed{\varepsilon}(p q)-\slashed{q}(\varepsilon p)] q_{\mu} \gamma_{5}\\
        \nonumber-&G_{1} \slashed{\varepsilon} q\{m+\slashed{p}\} q_{\mu} \gamma_{5}\\
        \nonumber+&\text { other structures with } \gamma^{\mu} \text { at the end}\\
        \nonumber+&\text{structures which are proportional to }(p+q)^{\mu}\Bigg\}
    \end{align}
%
%
Now let us turn our attention to the calculation of the correlation function from the QCD side. After applying Wick theorem the correlation function given in Eq.-(\ref{spin12bfcorr}) becomes:
%

    \begin{align}
        \nonumber\Pi_\mu^S(p) &= i\int d^4xe^{ipx}\dfrac{6}{\sqrt{6}}\bra{0}\left\{\pro{Q'}\gamma_5\Tr\Big[\gamma_\mu\pro{Q}\tilpro{q}\Big]\right.\\
        \nonumber&-\pro{Q}\tilpro{q}\gamma_\mu\pro{Q'}\gamma_5-\pro{q}\tilpro{Q}\gamma_\mu\pro{Q'}\gamma_5\\
        \nonumber&-\pro{Q'}\tilpro{q}\gamma_\mu\pro{Q}\gamma_5+\pro{Q}\gamma_5\Tr\Big[\tilpro{q}\gamma_\mu\pro{Q'}\Big]\\
        \nonumber&+\pro{q}\tilpro{Q'}\gamma_\mu\pro{Q}\gamma_5+\beta\pro{Q'}\Tr\Big[\gamma_\mu\pro{Q}\gamma_5\tilpro{q}\Big]\\
        \nonumber&-\beta\pro{Q}\gamma_5\tilpro{q}\gamma_\mu\pro{Q'}-\beta\pro{q}\gamma_5\tilpro{Q}\gamma_\mu\pro{Q'}\\
        \nonumber&-\beta\pro{Q'}\gamma_5\tilpro{q}\gamma_\mu\pro{Q}+\beta\pro{Q}\Tr\Big[\tilpro{q}\gamma_\mu\pro{Q'}\gamma_5\Big]\\
        \nonumber&\left.+\beta\pro{q}\gamma_5\tilpro{Q'}\gamma_\mu\pro{Q}\right\}\ket{0}_F\\
        \label{transition_after_wick}\Pi_\mu^A(p) &= i\int d^4xe^{ipx}\dfrac{6}{3\sqrt{2}}\bra{0}\left\{2\pro{Q'}\tilpro{Q}\gm{\mu}\pro{q}\gf-\beta\pro{q}\gf\tilpro{Q'}\gm{\mu}\pro{Q}\right.\\
        \nonumber&+2\pro{Q}\tilpro{Q'}\gm{\mu}\pro{q}\gf +2\pro{q}\gf\Tr\Big[\tilpro{Q}\gm{\mu}\pro{Q'}\Big]\\
        \nonumber&-\pro{Q'}\gf\Tr\Big[\pro{Q}\gm{\mu}\pro{q}\Big] -\pro{Q}\tilpro{q}\gm{\mu}\pro{Q'}\gf\\
        \nonumber&-\pro{q}\tilpro{Q}\gm{\mu}\pro{Q'}\gf +\pro{Q'}\tilpro{q}\gm{\mu}\pro{Q}\gf\\
        \nonumber&+\pro{Q}\gf\Tr\Big[\tilpro{Q'}\gm{\mu}\pro{q}\Big] -\pro{q}\tilpro{Q'}\gm{\mu}\pro{Q}\gf\\
        \nonumber&+2\beta\pro{Q'}\gf\tilpro{Q}\gm{\mu}\pro{q} -2\beta\pro{Q}\gf\tilpro{Q'}\gm{\mu}\pro{q}\\
        \nonumber&+2\beta\pro{q}\Tr\Big[\gm{\mu}\pro{Q'}\gf\tilpro{Q}\Big] -\beta\pro{Q'}\Tr\Big[\gf\tilpro{q}\gm{\mu}\pro{Q}\Big]\\
        \nonumber&-\beta\pro{Q}\gf\tilpro{q}\gm{\mu}\pro{Q'} -\beta\pro{q}\gf\tilpro{Q}\gm{\mu}\pro{Q'}\\
        \nonumber&\left.+\beta\pro{Q'}\gf\tilpro{q}\gm{\mu}\pro{Q} +\beta\pro{Q}\Tr\Big[\gf\tilpro{Q'}\gm{\mu}\pro{q}\Big]\right\}\ket{0}_F
    \end{align}
%
where the superscripts $S(A)$ denotes when symmetric(anti-symmetric) interpolating current has been used. In Eq-(\ref{transition_after_wick}), $S_Q$ and $S_q$ are heavy and light quark propagators. Their expressions in the presence of gluonic and electromagnetic background fields are
%
\begin{equation}
    \begin{aligned}
        S_q(x) &= \dfrac{i\slashed{x}}{2\pi^2x^4}-\dfrac{i}{16\pi^2x^2}\int_0^1du\Bigg\{\overline{u}\slashed{x}\sigma_{\alpha\beta}+u\sigma_{\alpha\beta}\slashed{x}\Bigg\}\Bigg\{gG^{\alpha\beta}(ux)+e_qF^{\alpha\beta}(ux)\Bigg\}\\
        S_Q(x) &= \dfrac{m_Q^2}{4\pi^2}\left[\dfrac{K_1\left(m_Q\sqrt{-x^2}\right)}{\sqrt{-x^2}}+\dfrac{i\slashed{x}}{\left(\sqrt{-x^2}\right)^2}K_2\left(m_Q\sqrt{-x^2}\right)\right]\\
        &-\dfrac{im_Q}{16\pi^2}\int_0^1 du\Bigg\{gG^{\alpha\beta}(ux) + e_QF^{\alpha\beta}(ux)\Bigg\}\\
        &\Bigg\{\sigma_{\alpha\beta}K_0\left(m_Q\sqrt{-x^2}\right) + \left(\left(\overline{u}\slashed{x}\sigma_{\alpha\beta}+u\sigma_{\alpha\beta}\slashed{x}\right)\dfrac{K_1\left(m_Q\sqrt{-x^2}\right)}{\sqrt{-x^2}}\right)\Bigg\}
    \end{aligned}
    \label{propagators}
\end{equation}
Where $K_i$ are the modified Bessel functions of the second kind, $G_{\mu\nu}$ and $F_{\mu\nu}$ are the background gluon and electromagnetic field strength tensors respectively.

Calculation of the correlation function in the external field involves perturbative and non-perturbative parts, i.e. the photon interacts with the quarks perturbatively or non-perturbatively. To calculate the perturbative part, when the propagators in Eq.-(\ref{propagators}) is placed in Eq.-(\ref{transition_after_wick}), the terms linear in $F_{\mu \nu}$ are selected. For the calculation of the non-perturbative part, it is enough to replace the light quark propagator that emits the photon with
\begin{equation}
    S_q \longrightarrow -\dfrac{1}{4}\Gamma_j\overline{q}\Gamma_j q
\end{equation}
where $\Gamma_j = \left\{\mathbf{1},\gamma_\alpha, i\gamma_5\gamma_\alpha, \gamma_5, \dfrac{1}{\sqrt{2}}\sigma_{\alpha\beta}\right\}$. In this case, there are two and three particle matrix elements: $\mel{0}{\overline{q}\Gamma_j q}{0}_F$, $\mel{0}{\overline{q}\Gamma G_{\mu\nu}q}{0}_F$ and $\mel{0}{\overline{q}\Gamma F_{\mu\nu}q}{0}_F$ which are parametrized in terms of photon DAs and describes the interaction of photons with quark fields at large distance.

We see that the hadronic and QCD sides of the correlation function contain many structures. Among all structures, the structures $\slashed{\varepsilon}\slashed{p}\gf q_\mu$ and $\slashed{q}\slashed{p}\gf\left(\varepsilon p\right)q_\mu$ exhibit the best convergence, for this reason, we choose them for determination of $G_1(0)$ and $G_2(0)$ respectively. equating the coefficients of the aforementioned structures on both sides, we get the sum rules for $G_1(0)$ and $G_2(0)$. In the final step, we perform Borel transformation over variables $-p^2$ and $-(p+q)^2$ in order to suppress the higher states and continuum contributions and enhance ground states to obtain
%
\begin{equation}
    \begin{aligned}
        \left(m_2+m_1\right)G_1(0)\lambda_1\lambda_2e^{-\nicefrac{m_2^2}{M_1^2}}e^{-\nicefrac{m_1^2}{M_2^2}}  +\cdots &= \int ds_1 \int ds_2 e^{-\nicefrac{s_1}{M_1^2}-\nicefrac{s_2}{M_2^2}}\rho_1(s_1,s_2)\\
        G_2(0)\lambda_1\lambda_2e^{-\nicefrac{m_2^2}{M_1^2}}e^{-\nicefrac{m_1^2}{M_2^2}} +\cdots&= \int ds_1 \int ds_2 e^{-\nicefrac{s_1}{M_1^2}-\nicefrac{s_2}{M_2^2}}\rho_2(s_1,s_2)\\
    \end{aligned}
    \label{correlated}
\end{equation}
%
where $\lambda_1$ and $\lambda_2$ are the residues of the spin-$\nicefrac{3}{2}$ and spin-$\nicefrac{1}{2}$ baryons respectively, and $\cdots$ denote the contribution from higher states and the continuum. To obtain the sum rules for the form factors $G_1$ and $G_2$, the contributions of the higher states and the continuum are subtracted using quark hadron duality: 
\begin{equation}
    \rho(s_1,s_2) \simeq \rho^{OPE}(s_1,s_2) ~~\mbox{if}~~ (s_1,s_2) \not\in {\cal D}
\end{equation}
where $\cal D$ is a domain in the $(s_1,s_2)$ plane.
Typically, the domain $\cal D$ is a rectangular region defined by $s_1<s_{10}$ and $s_2<s_{20}$ for some constants $s_{10}$ and $s_{20}$, or a triangular region. In this work, for its simplicity, continuum subtraction is carried out
by choosing $\cal D$ as the region defined as $s \equiv s_1 u_0 + s_2 \bar u_0 < s_0$ where $u_0\equiv \frac{M_2^2}{M_1^2+M_2^2}$ and $\bar u_0 = 1 - u_0$. Introducing a second variable $u=\frac{s_1 u_0}{s}$, the integral in the $(s_1,s_2)$ plane can be written as:
\begin{equation}
    \int ds_1 \int ds_2 e^{-s_1/M_1^2-s_2/M_2^2} \rho(s_1,s_2) = \int ds e^{-\frac{s}{M^2}} \rho(s)
\end{equation}
where 
\begin{equation}
    \rho(s) = \frac{s}{u_0 \bar u_0} \int_0^1 du \rho\left( s \frac{u}{u_0}, s \frac{\bar u}{\bar u_0} \right) 
\end{equation}
%
In the problem under study, the masses of the initial and final state baryons are very close to each other, hence we can set $M_1^2 = M_2^2 = 2M^2$ leading to $u_0=\nicefrac{1}{2}$. 
%
Using quark hadron duality, sum rules for the form factors take the form:
%
\begin{equation}
    \begin{aligned}
        \left(m_1+m_2\right)G_1(0)\lambda_1\lambda_2e^{-\nicefrac{m^2}{M^2}} &= \int_0^{s_0} ds e^{-\nicefrac{s}{M^2}}\rho_1(s) \\
        G_2(0)\lambda_1\lambda_2e^{-\nicefrac{m^2}{M^2}} &= \int_0^{s_0} ds e^{-\nicefrac{s}{M^2}}\rho_2(s)
    \end{aligned}
    \label{correlated_with_s}
\end{equation}
%
where $m^2=(m_1^2+m_2^2)/2$ and the expressions of spectral densities $\rho_1(s)$ and $\rho_2(s)$ are presented in Appendix-\ref{appndxb}. Note that although the limits of the $s$ integral is written from $s=0$ upto $s=s_0$, $\rho_i(s)=0$ ($i=1$ or $2$) for $s < (m_Q+m_{Q'})^2$.
%
From Eq-(\ref{correlated_with_s}) it follows that, for determination of $G_1(0)$ and $G_2(0)$ the residues $\lambda_1$ and $\lambda_2$ are needed. Spin-$\nicefrac{1}{2}$ residues are calculated in \cite{AlievOctet}, while the values for the spin-$\nicefrac{3}{2}$ residues are calculated in \cite{AlievDecuplet}.

